\section{Population Properties of Mergers Containing Neutron Stars}\label{sec:ns}

In O4a, GW230529~\citep{LIGOScientific:2024elc} is the only \ac{NS}-containing event identified with a ${\rm \ac{FAR} < 0.25\,{yr}^{-1}}$. We do not include the \ac{NSBH} candidate GW230518 identified during the \ac{ER} preceding \ac{O4a}. Additionally, no coincident \ac{EM} counterparts were identified for triggers (i.e., preliminary candidate GW signals flagged by the detection pipelines when their ranking statistic exceeded the alert threshold) during O4a based on follow-up efforts conducted by \ac{EM} telescopes. Thus, conclusions presented in previous papers~\citep{KAGRA:2021duu, LIGOScientific:2024elc} about the population properties of \ac{BNS} and \ac{NSBH} mergers remain largely unchanged.

We adopt models and methods consistent with those used in \gwtcthree~\citep{KAGRA:2021duu}. Specifically, we adopt two different models for the \ac{NS} mass distribution, one in which masses are assumed to be Gaussian distributed, and another in which they are assumed to follow a power law~\citep{KAGRA:2021duu}. These are referred to as the \textsc{Peak} model and the \textsc{Power} model respectively. We assume that the redshift evolution of the merger rate is fixed, and that the spins are distributed following the \ac{PE} prior~\citep{GWTC:Methods}. A uniform prior is used for hyperparameters of the population, with the condition \(m_\mathrm{min}\leq \mu \leq m_\mathrm{max}\) (where \(\mu\) is the mean of the Gaussian bump in the \textsc{Peak} model), assuming \(m_\mathrm{max}\) does not exceed \(M_\mathrm{max,TOV}\) (the maximum permissible \ac{NS} mass as expected from the TOV limit), as detailed in Appendix~\ref{appendix:parametric_models_summary}. 

When assumed to follow the \textsc{Power} model, \ac{NS} masses favor a mass distribution with a power-law slope constrained to $\alpha=\CIPlusMinus{\powersodapop[woGW190814][param][alpha]}$. The inference from the \textsc{Peak} model is broad, with largely unconstrained peak width $\sigma=\CIPlusMinus{\peakcutsodapop[woGW190814][param][sigma]}\,\Msun$ and location $\mu=\CIPlusMinus{\peakcutsodapop[woGW190814][param][mu]}\,\Msun$. While these results may hint at a peak emerging near \(1.4\,\Msun\), it is much broader than the relatively sharp peak in the mass distribution of Galactic \ac{NS} systems~\citep{Farrow:2019xnc,2024OJAp....7E..58E}. Our inferred \ac{NS} mass distribution remains broad, with greater support for high-mass \acp{NS}.

As no new \acp{NSBH} beyond GW230529 were confidently observed in O4a, the population-level results in \citet{LIGOScientific:2024elc} produced using the \textsc{NSBHPop} model~\citep{Biscoveanu:2022iue} remain unchanged. Specifically, our inferred minimum \ac{BH} mass in \ac{NSBH} systems remains {$3.4^{+1.0}_{-1.2}\,\Msun$}~\citep{LIGOScientific:2024elc}. 

